\section{Operating Systems Required or Recommended}\label{sec:os}

The sources do not say why Google chose Linux. They do show which properties of Linux matter for SRE work, and we describe those.

\paragraph{Production servers.}
The Linux kernel underpins much of Google's production infrastructure \cite{bsrs}. Google has an internal Linux distribution. It first installed every datacenter machine with a ``golden'' base image and later moved to separate release units for each software package \cite{bsrs}. A 2013 paper by a Google engineer describes a live upgrade of Google's servers from a Red Hat 7.1 image to a Debian Testing based distribution built from source, without a long ``flag day'', with applications running in a different partition from the base distribution \cite{usenix-conference-lisa13-technical-sessions-presentation-merlin}. Resource allocation in the datacenters is handled by Borg, which the SRE book calls Google's cluster operating system \cite{sre-book}. The Workbook says Borg is Google's internal container management system and runs huge numbers of Linux containers \cite{workbook}.

\paragraph{Workstations.}
A 2022 Google Cloud Blog post says Google operates many OS platforms for its employees, including a Linux system, across hundreds of thousands of devices \cite{google-cloud-blog}. Its internal Linux distribution, Goobuntu, was for a long time based on Ubuntu LTS. In 2018 Google completed a move to gLinux Rodete, short for Rolling Debian Testing, because Debian offered a smooth in-place migration. Each release is guided to the fleet with SRE practices such as incremental canarying, and the rolling schedule lets Google patch security holes on the whole fleet quickly \cite{google-cloud-blog}. No source says which OS Google SREs use on their own workstations. Our recommendation, which is an assumption, is to prepare to work every day in a Linux shell, both on a workstation and on remote servers.

\paragraph{Kernel updates and logs.}
Google regularly pushes new kernels to its whole fleet, with a target of less than 30 days \cite{bsrs}. It uses ksplice, a runtime kernel patch that avoids a reboot, and SREs applied one to production systems for two 2018 kernel vulnerabilities \cite{bsrs}. The Linux kernel documentation describes livepatching the same way \cite{kernel-docs}. For investigations, Linux provides syslog and auditd logs \cite{bsrs}.

\paragraph{Isolation and resource limits.}
Containers share the OS kernel \cite{cloud-google-learn-what-are-containers}. Linux cgroups organize processes in a hierarchy and distribute system resources along it, and Google contributed cgroups to the Linux kernel \cite{kernel-docs,cloud-google-learn-what-are-containers}. When a cgroup reaches its \texttt{memory.max} limit and the memory cannot be reduced, the kernel's out-of-memory killer runs inside that cgroup \cite{kernel-docs}. Google has many out-of-memory conditions every day \cite{bsrs}, and Borg kills and restarts a task that uses more resources than it requested \cite{sre-book}. A server can also run out of file descriptors or process IDs \cite{sre-book}. For untrusted code, Google added ptrace sandboxing to App Engine because the Linux kernel exposed a large attack surface \cite{bsrs}, and its security overview lists Linux user separation and an application kernel such as gVisor among its isolation techniques \cite{docs-cloud-google-docs-security-infrastructure-design}.
